\item El peso de un bloque rectangular de material de baja densidad es $15.0 \text{ N}$. Con una cuerda delgada, el centro de la cara inferior horizontal del bloque se amarra al fondo de un vaso de precipitados parcialmente lleno con agua. Cuando $25.0\%$ del volumen del bloque está sumergido, la tensión en la cuerda es $10.0 \text{ N}$. (a) Encuentre la fuerza de flotación sobre el bloque. (b) Ahora al vaso de precipitados se le agrega sin interrupción aceite de $800 \text{ kg/m}^3$ de densidad, lo que forma una capa sobre el agua y rodea al bloque. El aceite ejerce fuerzas sobre cada una de las cuatro paredes laterales del bloque que el aceite toca. ¿Cuáles son las direcciones de dichas fuerzas? (c) ¿Qué le ocurre a la tensión en la cuerda conforme se agrega el aceite? Explique cómo el aceite tiene este efecto sobre la tensión de la cuerda. (d) La cuerda se rompe cuando su tensión alcanza $60.0 \text{ N}$. En este momento, $25.0\%$ del volumen del bloque todavía está bajo la línea del agua. ¿Qué fracción adicional del volumen del bloque continúa por debajo de la superficie superior del aceite?
